\chapter{Jak zaczac prace z MoonChunk}

\section{Dwa poziomy pracy: uzytkownik jezyka i autor jezyka}
To bardzo wazne rozroznienie. W repozytorium MoonChunk mozna zobaczyc wiele polecen
zwiazanych z budowaniem projektu, generowaniem parsera czy uruchamianiem testow.
Sa one potrzebne glownie wtedy, gdy rozwijasz sam \emph{MoonChunk jako projekt}.

Ten podrecznik ma inny cel. Interesuje nas osoba, ktora chce \important{uzywac jezyka},
a niekoniecznie rozwijac jego implementacje. Z punktu widzenia takiego uzytkownika
najwazniejsze jest co innego:

\begin{itemize}
  \item jak napisac plik \mc{.mncnk},
  \item jak uruchomic jego wykonanie,
  \item gdzie pojawia sie wynik,
  \item jak odczytywac bledy,
  \item jak organizowac wiekszy projekt.
\end{itemize}

\section{Najprostszy model uzycia}
Obecny model pracy z MoonChunk najlatwiej opisac tak:

\begin{enumerate}
  \item instalujesz pakiet \mc{moonchunk} w projekcie Node.js,
  \item tworzysz swoje pliki \mc{.mncnk},
  \item uruchamiasz prosty skrypt JavaScript albo TypeScript,
  \item skrypt wywoluje funkcje \mc{executeMoonChunkFile(...)},
  \item MoonChunk generuje pliki HTML.
\end{enumerate}

To podejscie ma jedna duza zalete: nie musisz znac wewnetrznej struktury repozytorium,
a jednoczesnie masz pelna kontrole nad tym, kiedy i z jakiego pliku rozpoczynasz
budowanie strony.

\section{Instalacja w projekcie uzytkownika}
W zwyklym projekcie Node.js mozesz zainstalowac MoonChunk z rejestru npm.
Przykladowo:

\begin{lstlisting}[style=basecode,caption={Instalacja pakietu MoonChunk w projekcie uzytkownika}]
npm install moonchunk
\end{lstlisting}

Jesli korzystasz z Yarn:

\begin{lstlisting}[style=basecode,caption={Instalacja pakietu przez Yarn}]
yarn add moonchunk
\end{lstlisting}

Warto od razu przygotowac prosty katalog projektu, na przyklad taki:

\begin{lstlisting}[style=basecode,caption={Przykladowy uklad plikow projektu}]
my-site/
  content/
    site.mncnk
    posts.json
  run.js
  package.json
\end{lstlisting}

\section{Minimalny skrypt uruchamiajacy}
Poniewaz najwazniejszy interfejs dla uzytkownika jest obecnie programistyczny,
najprosciej utworzyc maly plik \mc{run.js}:

\begin{lstlisting}[style=basecode,caption={Najprostszy runner MoonChunk w JavaScript}]
const { executeMoonChunkFile } = require("moonchunk");

const result = executeMoonChunkFile("./content/site.mncnk");

if (!result.ok) {
  console.error("Wystapil blad:");
  console.error(result.diagnostics);
  process.exit(1);
}

console.log("Wygenerowane pliki:");
console.log(result.generatedFiles || []);
\end{lstlisting}

Ten skrypt jest prosty, ale w zupelnosci wystarczy do codziennego uzycia. Nie musisz
wiedziec nic o ANTLR, parserze ani o wewnetrznej architekturze projektu.

\section{Co robi \mc{executeMoonChunkFile(...)}}
Funkcja \mc{executeMoonChunkFile(...)} bierze sciezke do pliku \mc{.mncnk}, czyta
jego zawartosc, uruchamia parser i interpreter, a nastepnie zwraca obiekt z wynikiem.
Najwazniejsze pola tego obiektu to:

\begin{itemize}
  \item \mc{ok} -- informacja, czy wykonanie zakonczylo sie sukcesem,
  \item \mc{generatedFiles} -- lista wygenerowanych plikow HTML,
  \item \mc{diagnostics} -- lista komunikatow diagnostycznych, jesli wystapily bledy,
  \item \mc{output} -- dodatkowy log zwiazany z wykonaniem.
\end{itemize}

Dla zwyklego uzytkownika najczesciej wystarczy sprawdzenie pola \mc{ok}, a potem
odczytanie \mc{generatedFiles} albo \mc{diagnostics}.

\section{Dlaczego warto miec prosty runner}
Mozna zapytac: po co mi dodatkowy plik \mc{run.js}? Powod jest praktyczny.
Dzieki niemu w jednym miejscu decydujesz:

\begin{itemize}
  \item jaki plik \mc{.mncnk} jest punktem wejscia,
  \item co zrobic po sukcesie,
  \item jak wyswietlac bledy,
  \item czy chcesz generowac kilka projektow jeden po drugim.
\end{itemize}

To bardzo wygodny wzorzec, zwlaszcza gdy projekt rośnie.

\section{Pierwsze uruchomienie z przygotowanym plikiem}
Kiedy masz juz skrypt uruchamiajacy oraz plik \mc{site.mncnk}, wystarczy wywolac:

\begin{lstlisting}[style=basecode,caption={Uruchomienie generatora}]
node run.js
\end{lstlisting}

Jesli wszystko jest poprawne, MoonChunk wygeneruje gotowe pliki HTML. Gdzie trafi wynik?
To zalezy od tego, jaka sciezke podasz w instrukcji \mc{output}. O niej powiemy szerzej
w kolejnych rozdzialach.

\section{Najczestsze pytania poczatkujacych}
\subsection*{Czy musze znac JavaScript, zeby uzywac MoonChunk}
Nie w pelnym zakresie. Wystarczy Ci bardzo prosty skrypt startowy, ktory mozna raz
przygotowac i potem prawie do niego nie wracac. Sama praca nad trescia i logika stron
odbywa sie juz w MoonChunk.

\subsection*{Czy MoonChunk dziala tylko z Node.js}
Obecny sposob uzycia jest oparty o Node.js, bo wlasnie tam wykonywany jest pakiet.
Dla osoby korzystajacej z jezyka oznacza to po prostu tyle, ze potrzebuje srodowiska,
ktore uruchomi skrypt startowy.

\subsection*{Czy jedna strona to jeden plik}\
Niekoniecznie. Mozesz miec jeden plik z wieloma \mc{chunk}-ami, wiele plikow z importami,
a nawet wiele punktow wejscia przez \mc{moon(...)}. Oznacza to, ze strukture projektu
mozesz dopasowac do swojej skali pracy.

\section{Nastawienie, ktore pomaga}
MoonChunk najlatwiej poznawac malymi krokami. Dobra kolejnosc wyglada tak:

\begin{enumerate}
  \item najpierw zbuduj jedna, bardzo prosta strone,
  \item potem dodaj zmienne,
  \item nastepnie dodaj warunek i petle,
  \item dopiero potem przejdz do funkcji, importow i danych z JSON.
\end{enumerate}

W kolejnym rozdziale przejdziemy przez pierwszy kompletny program, zeby zobaczyc,
jak te wszystkie elementy lacza sie w jedna calosc.
